\section{Conclusiones}

La estructuración constituye una etapa fundamental del diseño sísmico, pues 
las decisiones tempranas sobre disposición y dimensionamiento de muros, vigas 
y losas condicionan directamente el comportamiento del edificio frente a 
sismos. Una estructuración adecuada garantiza continuidad en el flujo de 
fuerzas desde los niveles superiores hasta las fundaciones, evita 
concentraciones de esfuerzos, y permite anticipar el desempeño sísmico 
esperado antes de realizar un análisis dinámico completo. En edificios de 
mediana y gran altura como el presente, esta etapa resulta especialmente 
crítica, ya que errores en la concepción estructural son difíciles y costosos 
de corregir en etapas posteriores.

El sistema estructural adoptado corresponde a \textbf{muros de corte de 
hormigón armado}, lo que lo clasifica como un sistema de \textbf{muros 
estructurales} y no de marcos rígidos. Esta elección es consistente con 
la práctica habitual en edificios habitacionales de altura en Chile, donde 
la alta sismicidad exige sistemas de gran rigidez lateral. Los marcos 
rígidos, si bien son eficientes en zonas de menor sismicidad, presentan 
mayores desplazamientos laterales y requieren secciones considerablemente 
mayores para controlar derivas en zona sísmica~3. El sistema de muros 
provee la rigidez necesaria con un uso eficiente del hormigón, y es 
además compatible con la distribución arquitectónica del edificio, donde 
los muros divisorios entre departamentos cumplen simultáneamente función 
estructural.

Respecto al espesor de losas, se obtuvo un espesor de \textbf{10~cm} para 
todos los pisos habitacionales (pisos 1 al~21), valor que resultó gobernado 
por losas de relación de aspecto cercana a~1,5 bajo condiciones de apoyo 
Tipo~6 (cuatro bordes empotrados). Este espesor es razonable para las luces 
consideradas (menores a 6,25~m) y coherente con los rangos típicos de losas 
de edificios habitacionales en Chile. En el subterráneo, la losa~036 —con 
condición de apoyo Tipo~2b, que implica menor restricción en sus bordes— 
elevó el espesor determinante a 14,30~cm, adoptándose \textbf{15~cm} para 
ese nivel. La diferencia entre ambos espesores refleja directamente el efecto 
de las condiciones de borde sobre la esbeltez requerida: a menor 
empotramiento, mayor espesor necesario para controlar deflexiones.

En cuanto a la variación de espesor de muros en la altura, el criterio 
general indica que los espesores deben disminuir gradualmente hacia los 
niveles superiores, donde las solicitaciones sísmicas acumuladas son 
menores. Sin embargo, en el presente proyecto el índice de densidad de 
muros resultante ya en el nivel base se aproxima al límite inferior 
recomendado del 2{,}5\%, alcanzando valores de 2{,}7\% en dirección~$X$ 
para el piso tipo. Esta condición no permite reducir el espesor sin 
comprometer la densidad mínima exigida, por lo que se mantiene un espesor 
\textbf{uniforme de 25~cm en toda la altura del edificio}. Esta decisión 
es conservadora y queda sujeta a revisión en la etapa de diseño definitivo, 
una vez disponible el análisis modal espectral completo que permita 
evaluar con mayor precisión la demanda en cada nivel.